\section{相关工作}
\label{sec:related_work}

\paragraph{动态操作。}补偿感知与执行之间的时间差是一个经典问题，从带纯滞后的控制 \cite{smith1957}，到用卡尔曼类滤波器跟踪目标并规划拦截的抓取系统 \cite{barshalom2001,allen1993,baeuml2010,marturi2019dynamic}；这些系统精确，但针对每种物体和任务专门设计。近来，学习型策略开始进入运动场景 \cite{dynamicvla}。PUMA 读取历史光流，并以未来物体特征作为监督 \cite{domino}；DynamicWAM 以光流和运动学描述子为 WAM 提供条件，并以实时分块方式执行 \cite{dynamicwam}。两者都从图像空间线索中推断过去的运动，而 \ours 以目标的度量 SE(3) 位姿历史为输入，只观测当前一帧 RGB-D。另一些动态策略预测显著 token \cite{ahead}、未来图像 \cite{dbctfp}、紧凑的未来编码 \cite{foretime} 或 3D 场景动力学 \cite{physmani}，并在测试时以这些预测为条件生成动作。DynaWM 针对动态操作调整基础 VLA 的动作块 \cite{dynawm}，DynamicManip 从单条静态示范合成动态示范 \cite{dynamicmanip}。执行调度方案通过让推理与执行重叠 \cite{dynamicvla,smolvla}、保持相邻动作块一致 \cite{rtc,bid}，或以执行时刻为条件 \cite{vlash}，来减小或掩盖延迟。它们决定的是动作何时计算、何时执行，而不是策略表示什么。DOMINO \cite{domino} 与 PhysMani-Bench \cite{physmani} 提供了多样的动态操作任务，但它们公开的默认评测流程只在策略返回预测之后才推进仿真。\bench 则在推理期间推进场景，并让每个方法承担其实测延迟。

\paragraph{WAM 中的预测监督。}世界—动作模型（WAM）在同一个网络中生成未来视频与动作 \cite{uwm,uva,dreamzero,openwam_alpha}，其基础是把视频预测作为策略目标 \cite{gr1,seer}。由于生成视频代价高，近期工作改变了预测的内容，例如动态区域、深度与语义特征 \cite{dreamvla}，与 RGB 并列的光流和几何 \cite{dreamwam}，物体掩码或物体槽 \cite{maskwam,oawam}，以及紧凑隐变量 \cite{flare,wog}；它们也在追问测试时是否需要预测。Fast-WAM 在不做测试时想象的情况下保留了视频建模的大部分收益 \cite{fastwam}，PUMA 与 SG-WAM 则在训练后去掉预测器 \cite{domino,sgwam}。\ours 沿用这种只在训练中使用预测的方式，但主要预测的是度量 4D 状态，而不是像素、特征或隐变量，因此目标未来的位姿得到显式监督，而不是停留在隐式。

\paragraph{以物体为中心与基于几何的运动表示。}度量 3D 输入使策略更精确、样本效率更高 \cite{peract,dp3}。运动已被表示为 2D 点轨迹 \cite{atm,track2act,im2flow2act}、3D 流 \cite{generalflow,flowdreamer,pointworld,physmani}，以及物体关键点或 SE(3) 轨迹 \cite{kpam,spot,objectforesight}；刚体分解也早已用于学习场景动力学 \cite{se3nets,se3posenets}。这些工作用同一种表示描述场景：要么处处是稠密的流或轨迹，要么每个物体一个位姿。以目标为中心的预测器（包括预测物体层面未来的动态 VLA）只建模目标的未来 \cite{domino,spot}，把阻挡或推动目标的物体排除在外。DOCIR 把机器人、目标与障碍物分开，但只把这种区分用作输入编码 \cite{docir}。\ours 按角色分配精度：粗粒度的稀疏 3D 流描述整个场景，接触依赖毫米级精度的目标另由其 SE(3) 历史与预测描述。
